\chapter{Installation}
\label{chap:installation}

\section{Requirements}

\actsim requires Python 3.8 or later. The package depends on a number of
common scientific Python libraries:

\begin{itemize}[leftmargin=*]
    \item \texttt{numpy} ($\geq 1.26$)
    \item \texttt{scipy} ($\geq 1.7$)
    \item \texttt{pandas} ($\geq 1.3$)
    \item \texttt{matplotlib} ($\geq 3.4$)
    \item \texttt{seaborn} ($\geq 0.13$)
    \item \texttt{actstats} ($\geq 0.1.1$)
\end{itemize}

These packages will be installed automatically when \actsim is installed
via \texttt{pip}.

\section{Standard Installation}

\actsim can be installed from the Python Package Index using \texttt{pip}:

\begin{lstlisting}[style=bashstyle]
pip install actsim
\end{lstlisting}

This will install the latest stable release together with all required
dependencies.

\section{Installation from Source}

To install the latest development version directly from the source
repository, clone the repository and install in editable mode:

\begin{lstlisting}[style=bashstyle]
git clone https://github.com/jzhng105/actsim.git
cd actsim
pip install -e .
\end{lstlisting}

Editable installation is convenient when modifying the package locally,
since changes to the source files take effect without reinstallation.

\section{Development Installation}

Contributors who plan to add features, fix bugs, or run the test suite
should install the additional development dependencies:

\begin{lstlisting}[style=bashstyle]
git clone https://github.com/jzhng105/actsim.git
cd actsim
pip install -e .[dev]
\end{lstlisting}

The development extras include testing utilities such as \texttt{pytest},
formatters, and other tools used in the contribution workflow.

\section{Recommended Virtual Environment}

It is recommended to install \actsim into an isolated virtual environment
to avoid conflicts with other Python projects:

\begin{lstlisting}[style=bashstyle]
python -m venv venv
source venv/bin/activate        # macOS / Linux
venv\Scripts\activate           # Windows
pip install actsim
\end{lstlisting}

\section{Verifying the Installation}

After installation, verify that \actsim can be imported and that the
default configuration loads successfully:

\begin{lstlisting}[style=pythonstyle]
from actsim import (
    DistributionFitter,
    StochasticSimulator,
    ClaimSimulator,
    load_config,
)

config = load_config()
print("Severity distributions:", config.distributions["severity"])
print("Frequency distributions:", config.distributions["frequency"])
print("Metrics:", config.metrics)
\end{lstlisting}

If the import statements succeed and the configuration prints the expected
lists of distributions and metrics, the installation is working correctly.
